% titel: Blatt 6, A1: Elektronegativität vs. Elektronenaffinität
% skript: Mortimer Kap. 8 (chemie: Atombindung, Elektronegativität)

Elektronegativität und Elektronenaffinität werden oft verwechselt. Beide
beschreiben, wie stark ein Atom Elektronen anzieht --- aber in verschiedenen
Situationen und mit unterschiedlicher Messbarkeit.

\begin{werkzeugkasten}[EN vs. EA]
\begin{itemize}
  \item \textbf{Elektronenaffinität (EA):} Energieänderung, wenn ein
        \textit{isoliertes, gasförmiges Atom} ein Elektron aufnimmt
        ($A(\mathrm g)+e^- \to A^-(\mathrm g)$). Sie ist klar definiert und
        experimentell messbar (\si{\kilo\joule\per\mol}).
  \item \textbf{Elektronegativität (EN):} relativer Zahlenwert, der angibt, wie
        stark ein \textit{in einer Bindung} gebundenes Atom die gemeinsamen
        Elektronen an sich zieht. Sie ist nicht direkt messbar, sondern nur über
        Skalen (z.\ B. Pauling) abgeschätzt.
  \item Trends: EN nimmt im PSE nach rechts und nach oben zu; F ist am
        elektronegativsten. EA ist am ausgeprägtesten bei Halogenen.
\end{itemize}
\textit{Skript:} Mortimer Kap. 8 (Elektronegativität, Elektronenaffinität).
\end{werkzeugkasten}

\begin{achtung}[Darauf achten]
\begin{itemize}
  \item EN ist eine Grosse für \textit{gebundene} Atome, EA für \textit{freie}
        Atome --- das ist der Kernunterschied.
  \item Nur die EA hat eine eindeutige, messbare Zahl; EN-Werte sind
        Konventionen.
  \item Für den Vergleich zweier Atome: EN wächst nach rechts/oben --- F ist
        das elektronegativste Element.
\end{itemize}
\end{achtung}
